% TOP_PROBLEM: {"id": 30006736, "title": "Prove the BKL locality conjecture in the general inhomogeneous context", "category_id": 16, "category": {"id": 16, "name": "physics", "display_name": "Mathematical Physics", "description": "Problems at the intersection of mathematics and physics.", "slug": "physics", "order_index": 16, "created_at": "2026-07-31T15:26:25.671Z"}, "status": "open"}
% FIRST_POSTED: 2026-09-25
% ENTRY_KIND: statement_only
% !TeX program = lualatex
% Definition and sources only; this file is not an LLM solution attempt.
\documentclass[11pt]{article}
\usepackage[margin=25mm]{geometry}
\usepackage{fontspec}
\setmainfont{Latin Modern Roman}
\usepackage{amsmath,amssymb,mathtools}
\usepackage{unicode-math}
\setmathfont{Latin Modern Math}
\usepackage{xurl,hyperref}
\hypersetup{colorlinks=true,urlcolor=blue}
\setcounter{secnumdepth}{0}
\setlength{\parindent}{0pt}
\setlength{\parskip}{5pt}
\setlength{\emergencystretch}{3em}
\begin{document}
\section*{\texorpdfstring{Prove the BKL locality conjecture in the general inhomogeneous context}{Prove the BKL locality conjecture in the general inhomogeneous context}}\label{30006736}
\subsection{Definitions and mathematical statement}
Near a spacelike singularity, a Kasner epoch has a schematic local metric
\[
 ds^2\simeq-dt^2+\sum_{i=1}^3t^{2p_i}(\omega^i)^2,\qquad
 \sum_i p_i=\sum_i p_i^2=1.
\]
Vacuum Bianchi IX (Mixmaster) behavior consists of successive anisotropic epochs and transitions in a spatially homogeneous model. BKL locality predicts that, for generic inhomogeneous collapse, time evolution asymptotically dominates coupling between distinct spatial points, with corresponding oscillatory local behavior. The catalogue does not specify a data topology or measure defining ``generic,'' a gauge, or a convergence mode. Consequently this is an asymptotic geometric program, not a fully quantified theorem determined by the single sentence alone.
\par\smallskip\textit{Formulation and concept references:} \href{https://link.springer.com/article/10.1007/s41114-025-00058-z}{[S1]}.

\subsection{Short English statement}
Near a generic gravitational singularity, does each spatial point evolve almost independently through oscillatory Mixmaster-like dynamics?
\subsection{Sources}
\begin{itemize}
\item[S1]Josu C. Aurrekoetxea, Katy Clough, and Eugene A. Lim, Cosmology using numerical relativity, Section 3.1 (2025).
\url{https://link.springer.com/article/10.1007/s41114-025-00058-z}.
\textit{Formulation source.}
\end{itemize}
\end{document}
